\section{Benchmark e Difesa dell\textquotesingle Architettura (AI
GENERATO)}\label{benchmark-e-difesa-dellarchitettura-ai-generato}

Spunti:

\begin{itemize}
\tightlist
\item
  perche\textquotesingle{} Debian 13?
\item
  perche\textquotesingle{} ZFS?
\item
  perche\textquotesingle{} OVN in verifica?
\item
  perche\textquotesingle{} Birdge L2 in produzione?
\item
  perche\textquotesingle{} Incus?
\end{itemize}

Codice:

\begin{itemize}
\tightlist
\item
  Motivazione per Async Context Manager
\item
  perche disaccoppiare la creazione dei comandi?
\item
  perche gestire le API in modo asincorno?
\item
  percher usare httpx.AsyncClient per IncusEngine?
\item
  perche BlueAgent gestisce i container e non mantiene lo stato delle
  challenge?
\item
  motivare le API rests
\item
  perche SQLite
\end{itemize}

\section{Critica e Limiti
Architetturali}\label{critica-e-limiti-architetturali}

Nonostante i vantaggi, l\textquotesingle architettura presenta
compromessi e vincoli strutturali:

\begin{itemize}
\tightlist
\item
  \textbf{UDP Path MTU Blackhole:} La combinazione di tunnel WireGuard
  (\(\text{MTU } 1420\)) e Geneve OVN (\(\text{MTU } 1442\),
  all\textquotesingle interno della \textbf{rete di verifica}) impone
  riduzioni di MTU. Mentre il TCP beneficia del MSS clamping via
  \passthrough{\lstinline!nftables!}, i pacchetti UDP di grandi
  dimensioni con flag \emph{Don\textquotesingle t Fragment}, potrebbero
  essere scartati dal kernel.
\item
  \textbf{Limite scalabilità a 254 team:}
  L\textquotesingle indirizzamento deterministico
  \passthrough{\lstinline!10.k.x.x!} vincola \(k\) a un singolo byte
  (\(0 \le k \le 254\)), quindi ad un massimo di \(\text{255}\) team
\item
  \textbf{Concorrenza e blocco multi-nodo:} Il meccanismo
  \passthrough{\lstinline!TwoTierLock!}
  (\passthrough{\lstinline!asyncio.Lock!} combinato con
  \passthrough{\lstinline!flock!}) è confinato al singolo nodo. Per
  scalare l\textquotesingle architettura su piu nodi bisogna introdurre
  l\textquotesingle uso di un gestore di lock distribuito come Redis.
\end{itemize}
